\section{Real-data admissibility study: testing the model class}
\label{sec:real-admissibility}
\subsection{Question, data reduction, and support}
The preceding synthetic fit starts with a known admissible field. Real
electrical samples need not satisfy that premise. This study asks whether
one scalar resistance adjustment can make the supplied terminal-current
maps compatible with both reflection and conservation. It is a stress test
of the assumed model class, not a demonstrated reconstruction of real
magnetic co-energy. No potential fit is performed on these real data.

Two supplied MATLAB v7.3 measurement files are processed through their
existing importer and stationary voltage inversion. One contains 2000-rpm
slices at rotor-temperature references 40, 60, and 80 $^\circ$C, with three
pole pairs and used $R=\SI{7.695}{\milli\ohm}$. The other contains 1000/3000-rpm
slices at a 70-$^\circ$C reference, four pole pairs, and used
$R=\SI{10}{\milli\ohm}$. They are not pooled as one machine calibration.
Neither a reference temperature nor a configured resistance verifies the
actual winding state. Requested speed sets $\omega_e=2\pi pn/60$; recorded
sample speeds do vary.

Signals are first averaged within operating-index groups. Requested d/q
currents rounded to 0.01 A define keys; repeated keys receive equal weight
across those groups. Measured currents and voltages are not rounded.
Reflection pairs use exactly mirrored requested keys, without filling
missing partners. Actual measured currents need not be exact mirrors.
For each pair let $r_d$ be the d half-difference, $r_q$ the q half-sum,
and $a_q=(|i_q^+|+|i_q^-|)/2$. The descriptive equivalent is
\begin{equation}
 \resistanceequivalent=-\omegae r_d/a_q,\qquad
 \symmetryresistance=R-\operatorname{median}_{\rm high}
                      (\resistanceequivalent).
 \label{eq:real-sym-equivalent}
\end{equation}
Division is accepted for $a_q>0.10a_{q,\max}$; the empirical high subset has
$a_q\ge0.85a_{q,\max}$. This is a conditional summary, not an independent
resistance measurement. Nonlinear parity leakage from imperfect measured
pair geometry remains possible.

The path study instead uses both independent componentwise linear
interpolants on a common Delaunay mesh in measured current coordinates.
Targets must exceed ten percent of each axis maximum in absolute magnitude;
both origin-to-target paths must lie within the convex hull. Hull acceptance
does not establish dense coverage: triangles can bridge unmeasured holes.
\Cref{tab:audit-support} exposes the distinct pair and rectangle support.
These are overlapping rectangles, not independent resistance experiments.
\begin{table}[tb]
\centering\small
% Generated by numerics/export_audit_tables.py; do not edit.
% audit_results.json SHA256 1523aedbba70a4364963d65cac529a522b4208336296fdc04ce6b2cdff27a765
% observable_checks.json SHA256 f1235618a0bf5fe74cc0a84385a49cb39d5da9a8db0b322ebd75022387a00d54
\begin{tabular}{rrrrr}
\toprule
$n$ & $T_{\rm rot,ref}$ & OP/keys & Pairs/high & Rectangles/skips \\
\midrule
2000 & 40 & 4220/4154 & 1172/79 & 2868/3 \\
2000 & 60 & 4220/4154 & 1172/80 & 2869/2 \\
2000 & 80 & 4220/4154 & 1172/79 & 2839/32 \\
1000 & 70 & 109/104 & 39/6 & 65/4 \\
3000 & 70 & 109/104 & 39/7 & 45/24 \\
\bottomrule
\end{tabular}

\caption{Unchanged audit support. Speed $n$ is in rpm and rotor reference in
$^\circ$C. Skips count hull failures after the axis threshold: 2871 candidate
targets per 2000-rpm slice and 69 per multi-speed slice. OP means
operating-index groups, not raw samples.}
\label{tab:audit-support}
\end{table}

\subsection{Oriented paths measure averaged circulation}
Write $x=i_d$, $y=i_q$. Let $\pathintegralA$ be the dq flux integral along
$(0,0)\to(x,0)\to(x,y)$ and $\pathintegralB$ along
$(0,0)\to(0,y)\to(x,y)$. These are reconstructed flux integrals; the notation
does not assume that a single-valued co-energy exists. Green's theorem with
oriented limits gives, in every quadrant,
\begin{align}
 \pathintegralA-\pathintegralB
 &=\int_0^x\int_0^y \fluxcurlresidual(\xi,\eta)\,d\eta\,d\xi
   =\operatorname{sgn}(xy)\iint_{\mathcal A_{\rm geom}}
     r_{\rm curl}\,dA,\label{eq:oriented-real-path}\\
 \pathresistanceerror(x,y)
 &=\frac{\omegae(W_A-W_B)}{2xy}
   =\frac{\omegae}{2|xy|}\iint_{\mathcal A_{\rm geom}}
     r_{\rm curl}\,dA.\label{eq:rectangle-equivalent}
\end{align}
Thus the endpoint indexes a rectangle. The result is an area-averaged
resistance-equivalent curl residual, not a local resistance at that endpoint.
The local derivative diagnostic would be
$\curlresistanceerror=\omegae\fluxcurlresidual/2$.
Opposite curl contributions can cancel within a rectangle; a finite set of
small loop integrals cannot prove local integrability everywhere.

For a sole resistance mismatch $\resistanceerror=R-R_{\rm true}$, voltage
inversion adds $\resistanceerror\rotationgenerator\current/\omegae$ to an
otherwise admissible field. Its curl is $2\resistanceerror/\omegae$ and
its two path contributions differ by $2\resistanceerror xy/\omegae$.
Only under that baseline does subtracting the equivalent recover
$R_{\rm true}$. More generally define the diagnostic
$\pathresistance=R-\pathresistanceerror$.

Why does this diagnostic remain unchanged when $R$ changes? At fixed
voltages, currents, speed and mesh,
\begin{equation}
 \flux_{R+h}^{\rm rec}=\flux_R^{\rm rec}
                  +\frac{h}{\omegae}\rotationgenerator\current.
\end{equation}
Linear barycentric interpolation reproduces this affine addition exactly;
trapez integration also integrates it exactly. Consequently
\begin{equation}
 \pathresistanceerror(R+h)=\pathresistanceerror(R)+h,
 \qquad
 \boxed{(R+h)-\Delta R_{\rm eq,path}(R+h)
              =R-\Delta R_{\rm eq,path}(R).}\label{eq:path-invariance}
\end{equation}
The 60-$^\circ$C calculation retains the pointwise value within
$1.39\cdot10^{-16}$ $\Omega$ after changing $R$ to the symmetry equivalent.
Additional tests use arbitrary shifts in all four quadrants.
This is a strong algebra/implementation check, not an independent calibration
of true winding resistance. Other circulating contributions survive it.

\subsection{The real discrepancy survives exact integration}
\Cref{tab:audit-equivalents} compares the five slices without treating the
different summaries as estimates of one proven parameter. At 2000 rpm and
60 $^\circ$C the high-current symmetry equivalent is
\SI{8.274010}{\milli\ohm}; the path median is
\SI{5.885553}{\milli\ohm} with raw MAD
\SI{1.728698}{\milli\ohm}. Exact integration of the same P1 field gives
\SI{5.885512}{\milli\ohm}. Here ``exact'' qualifies the interpolator's
integration, not its magnetic accuracy. MAD describes sampled spread and
is not a calibrated uncertainty interval.
\begin{table}[tb]
\centering\small
% Generated by numerics/export_audit_tables.py; do not edit.
% audit_results.json SHA256 1523aedbba70a4364963d65cac529a522b4208336296fdc04ce6b2cdff27a765
% observable_checks.json SHA256 f1235618a0bf5fe74cc0a84385a49cb39d5da9a8db0b322ebd75022387a00d54
\begin{tabular}{rrrrrrr}
\toprule
$n$ & $T_{\rm rot,ref}$ & $N_{\rm rect}$ & $R_{\rm eq,sym}$ & $R_{\rm eq,path}$ & MAD & P1 exact \\
\midrule
2000 & 40 & 2868 & 7.599140 & 5.007865 & 1.783325 & 5.007792 \\
2000 & 60 & 2869 & 8.274010 & 5.885553 & 1.728698 & 5.885512 \\
2000 & 80 & 2839 & 8.963406 & 6.488047 & 1.769623 & 6.487975 \\
1000 & 70 & 65 & 12.156625 & 5.960078 & 2.405193 & 5.960063 \\
3000 & 70 & 45 & 15.334054 & 3.809622 & 2.991105 & 3.808417 \\
\bottomrule
\end{tabular}

\caption{Conditional symmetry equivalent, median path equivalent, its raw
MAD, and exact-P1 path median from existing audit JSON. All resistance
columns are m$\Omega$; $n$ is rpm and rotor reference is $^\circ$C.
The symmetry summary uses high pairs; path statistics use all accepted
rectangles. Six decimal places document reproducibility, not accuracy.}
\label{tab:audit-equivalents}
\end{table}

Each path segment is split at mesh edges and its affine pieces integrated
exactly. Independently, triangle clipping against each rectangle integrates
the piecewise constant curl over intersection areas. On eight deterministic
real rectangles the maximum difference between these calculations is
$2.54\cdot10^{-14}$ Wb A. Across all 2869 accepted 60-$^\circ$C targets,
the original 500-point trapez rule differs from exact P1 integration by only
\SI{0.0000766}{\milli\ohm} RMS and
\SI{0.000628}{\milli\ohm} maximum absolute error. Quadrature therefore does
not explain the roughly \SI{2.388}{\milli\ohm} discrepancy. This rules out
one numerical explanation, not errors of interpolation or model reduction.

The 1000/3000-rpm path support is 65 versus 45 rectangles. Their medians
are not a pointwise speed test. Under a sole constant resistance mismatch,
the flux error scales as $1/\omega_e$, but a voltage error proportional to
$-\current$ yields the identical reconstruction fingerprint. Extra speeds
alone do not separate all unknown sources. The companion's direct comparison
on 36 common symmetry pairs remains a separate check.

\subsection{Improved high-current parity does not establish global consistency}
Change only the resistance used in reconstruction from
\SI{7.695}{\milli\ohm} to the high-current symmetry equivalent
\SI{8.274010}{\milli\ohm}. This is a hypothetical adjustment; the measured
signals are retained. \Cref{tab:audit-adjustment} shows strong improvements
within the selected 80 pairs. Across all 1172 pairs, however, the d and
vector residual RMS increase. The vector RMS is
$\sqrt{\operatorname{mean}(r_d^2+r_q^2)}$ on the stated pair set.
\begin{table}[tb]
\centering\small
% Generated by numerics/export_audit_tables.py; do not edit.
% audit_results.json SHA256 1523aedbba70a4364963d65cac529a522b4208336296fdc04ce6b2cdff27a765
% observable_checks.json SHA256 f1235618a0bf5fe74cc0a84385a49cb39d5da9a8db0b322ebd75022387a00d54
\begin{tabular}{lllrr}
\toprule
Support & Metric & Unit & Raw & Adjusted \\
\midrule
High (80 pairs) & $r_d$ RMS & mWb & 1.100092 & 0.219939 \\
High (80 pairs) & $r_q$ RMS & mWb & 0.487494 & 0.277063 \\
High (80 pairs) & Vector RMS & mWb & 1.203268 & 0.353747 \\
All (1172 pairs) & $r_d$ RMS & mWb & 1.746480 & 1.954456 \\
All (1172 pairs) & $r_q$ RMS & mWb & 1.308799 & 1.183835 \\
All (1172 pairs) & Vector RMS & mWb & 2.182464 & 2.285030 \\
2869 rectangles & Median $\rho_{\rm path}$ & \% & 3.488090 & 4.226703 \\
2869 rectangles & Median $\Delta R_{\rm eq,path}$ & m$\Omega$ & 1.809447 & 2.388457 \\
2869 rectangles & Median $R_{\rm eq,path}$ & m$\Omega$ & 5.885553 & 5.885553 \\
\bottomrule
\end{tabular}

\caption{2000 rpm/60 $^\circ$C: raw reconstruction versus a symmetry-based
resistance hypothesis on unchanged support. High and all-pair parity are
separate metrics; the last three rows concern flux path integrals. No
experimentally identified correction is implied.}
\label{tab:audit-adjustment}
\end{table}
The relative path difference is
\begin{equation}
 \relativepathdifference=
 \frac{|W_A-W_B|}{(|W_A|+|W_B|)/2},\label{eq:relative-real-path}
\end{equation}
where the denominator is nonzero. Its median rises from 3.488090 to
4.226703 percent. This is a descriptive comparison of reconstructed
integrals, not error against measured magnetic energy. The difference of
these medians is 0.738613 percentage points; it is distinct from the median
pointwise change (about 0.875 percentage points). The path equivalent itself
is invariant, as predicted, while the circulation equivalent and path
disagreement grow. The adjustment improves the selected parity channel;
it does not improve global reflection structure or establish conservation.

\subsection{P1 interpolation and differentiation set different limits}
Componentwise linear Delaunay interpolation is $C^0$, not $C^1$.
Its component gradients are constant within a triangle and jump at edges.
A common mesh does not enforce a symmetric Jacobian or turn it into the
Hessian of a common potential. At mesh vertices the classical derivative
is not unique. Since the field itself is continuous, no field-jump term
is needed in the weak curl, and the piecewise Green check remains valid.

The real 60-$^\circ$C mesh contains 8234 triangles. Its local
$\Delta R_{\rm eq,curl}$ has median about $-0.217$ m$\Omega$ and raw MAD
$4.869$ m$\Omega$, but unweighted RMS about $2984$ m$\Omega$ and extremes
about $-155108$ and $+128366$ m$\Omega$. Maximum edge-matrix condition number
is about 15606, and a nearly collinear boundary triangle has circumradius
about \SI{788738}{\ampere}. These peaks diagnose numerical fragility and
must not be read as identified physical resistance. A purely geometric
condition-number cutoff at 10 retains 7991 triangles and 99.208 percent
of hull area, reducing RMS to about $10.438$ m$\Omega$. It is a sensitivity
check, not a validated physical filter; full results remain available.
Local linear regressions with 12, 24, and 48 neighbors also depend on scale
and boundary geometry. Their different weighting does not reproduce the
distribution of overlapping rectangle averages.

An independent conservative saturated polynomial was sampled at the 4220
actual operating-point coordinates and passed through the unchanged key
reduction and componentwise P1 interpolation. For the audit's 64 deterministic
targets its interpolated path equivalent has RMS
\SI{0.013565}{\milli\ohm} although its analytic circulation vanishes to
numerical precision. Sampling directly at aggregated coordinates gives
essentially the same RMS; the OP-to-key averaging contribution is only
\SI{0.00000216}{\milli\ohm} RMS in this example. Linear conservative fields
are recovered to roundoff. A +2-m$\Omega$ injection produces the expected
pointwise increment, while coherent rotation preserves analytic integrability
and changes parity. Imperfect measured pair geometry also gives small
baseline parity leakage even for the conservative reference. These tests
separate numerical mechanisms; the small nonlinear example is not an upper
error bound for the unknown real field.

\subsection{What the recorded quantities actually establish}
The files do not provide a verifiable abc-to-dq transformation definition.
The median recorded phase-current RMS divided by
$\sqrt{i_d^2+i_q^2}$ is approximately 0.7075, 0.7077, and 0.7071 in the
three thermal slices, and 0.7085/0.7046 in the speed slices. This is consistent
with amplitude-invariant dq currents and $1/\sqrt2$, but does not verify
voltage scaling, power balance, harmonic content, or matched averaging.
The checked phase-current/power instrument channels and AC voltage RMS
channel are zero and provide no independent scaling evidence.

Accordingly $W_A,W_B$ are initially dq flux integrals in Wb A, not independently
measured absolute magnetic Joule values. The potential normalization defined
in \cref{sec:potential} remains a consistent convention. A constant factor
$\kappa$ multiplies physical integrals and their resistance denominator
together: $\Delta R=\omega_e\Delta W_{\rm phys}/(2\kappa xy)$.
A global $3/2$ changes neither path independence nor this discrepancy.
The measurements require verification of coordinates and voltage provenance
before the fit can be interpreted as a magnetic storage law. The reproducible
audit JSON/CSV, source and dataset hashes, and generated manuscript tables
retain this evidence separately from the chosen model class.
